% D04-009. AI-04 P0259--P0273.
% reports/phase4/kan_20260616_072711/interpretability/channel_dropout_probing.json
\section{Влияние групп входных признаков на обученную модель}
\label{sec:ch4_input_sensitivity}
\conceptsource{Статья 4. Программное зануление входов и градиентный анализ дополняют среднее качество. Эти проверки не являются обучением новых сетей или готовым испытанием реального отказа.}


\iffalse
(Примечание автора: Бегло просмотрел и даже не стал вычитывать. В исходных работах явно намного больше и понятнее, особенно больше картинок как раз, а у тебя прост слова и всё, без какого-то отображения практически. Я за то, чтобы максимально полно переносить статьи, пусть и перекомбинируя их и объединяя мысли) **[AA100: рассмотрено 27.10; ожидает проверки]**
(Примечание автора: Хм…. Ок, закрывай, если что-то нужно, укажу локально.)
\fi
\authorreviewnote{AA100}

Сопоставление представлений в разделе~\ref{sec:ch4_signal_representations}
показывает результат обучения разных моделей
на разных входных данных. Дополнительную информацию
даёт отключение части признаков уже обученной сети.
Такой опыт позволяет установить, какие группы
используются в сформированном решении.
Его постановка близка к отключению арифметических
ветвей из раздела~\ref{sec:ch4_physical_validation},
но изменяется вход модели, а не состав её внутренних ветвей.

\subsection{Порядок проверки}
\label{sec:ch4_input_sensitivity_protocol}
\conceptsource{Статья 4. Группы спектрального входа и фиксированная обученная модель, отдельно сохранённые симметричные составляющие.}


В исследовании детекции событий для Physical KAN-Transformer
последовательно занулялись выбранные группы
спектральных признаков. После каждого изменения
повторно оценивались выходные метки,
а коэффициенты сети оставались неизменными.
Рассматривались фазные токи и напряжения,
модули и углы, симметричные составляющие
и отдельные группы гармоник.

Вход включает 208 полярных координат отдельных
измерительных каналов и 12 координат симметричных
составляющих. При занулении выбранных фазных каналов
или гармоник их ранее рассчитанные симметричные
составляющие сохраняются. Поэтому такое воздействие
не равно удалению соответствующей информации
из всего тракта измерения.

Сохранённая проверка охватывает 14\,972 окна
при ограничении в 960 файлов RTDS.
Шаг окон здесь составляет один период,
порог двоичных выходов равен 0,5.
Исходный Macro-F1 составляет 0,9577.
Это отдельная проверка чувствительности,
а не повтор всех условий таблицы~\ref{tab:ch4_ozz_comparison}.
В частности, число и положение окон отличаются,
поэтому её исходная оценка не обязана совпадать
с оценкой общей сравнительной серии.

Зануление полярного признака имеет определённый смысл.
Нулевой модуль исключает соответствующий вклад,
а нулевой угол задаёт конкретное направление вектора.
Нуль не является универсальным обозначением
неизвестного значения для обеих координат.
Поэтому эксперимент прежде всего характеризует
изменение ответа при заданном воздействии на вход.

\subsection{Изменение качества при занулении признаков}
\label{sec:ch4_input_sensitivity_results}
\conceptsource{Статья 4. Влияние токов, напряжений, углов, модулей и групп гармоник; число гармоник не оптимизируется одним занулением.}


Выбранные результаты приведены
в таблице~\ref{tab:ch4_input_sensitivity}.
Все оценки относятся к одному обученному состоянию
и одной проверочной процедуре.

\begin{table}[htbp]
 \centering
 \small
 \caption{Влияние зануления групп признаков на Macro-F1
 обученного Physical KAN-Transformer}
 \label{tab:ch4_input_sensitivity}
 \begin{tabular}{p{9.4cm}r}
  \hline
  Вариант входа & Macro-F1 \\
  \hline
  Все признаки & 0,9577 \\
  Без фазных напряжений & 0,0160 \\
  Без фазных токов & 0,9537 \\
  Без модулей всех спектральных векторов & 0,0191 \\
  Без углов всех спектральных векторов & 0,0864 \\
  Без первой гармоники отдельных каналов & 0,1333 \\
  Без гармоник со второй по девятую отдельных каналов & 0,1992 \\
  Без субгармонических признаков & 0,8940 \\
  Только симметричные составляющие & 0,0056 \\
  \hline
 \end{tabular}
\end{table}

Наибольшее снижение среди физических групп каналов
наблюдается при отключении фазных напряжений.
Отключение фазных токов в этой проверке
влияет значительно слабее.
Полученное различие согласуется с существенной ролью
изменения фазных напряжений при ОЗЗ
в сети с изолированной нейтралью.
Оно относится к выбранной задаче и обученной модели.
В частности, его нельзя переносить на РНМ,
где совместные токовые и напряженческие величины
задают направление мощности.

Для полярного входа существенны обе координаты.
Зануление модулей или всех углов резко снижает качество.
Этот результат поддерживает рассмотрение
фазовых соотношений совместно с амплитудами,
но не выделяет единственный способ кодирования угла
для всех архитектур.

Отключение первой гармоники отдельных каналов
и отключение их гармоник со второй по девятую
также приводят к большим потерям.
При этом в обоих случаях остаётся значительная
часть остальных входных признаков.
Таким образом, обученная модель использует
как основную частоту, так и более высокие гармоники.
Слабее влияет совместное исключение
субгармонических признаков.

Эти опыты не устанавливают оптимальное число гармоник.
Модель после зануления части входа не обучается заново.
Обучение компактного варианта может иначе распределить
зависимости между оставшимися признаками.
По той же причине малое качество при сохранении
только симметричных составляющих не противоречит
результатам раздела~\ref{sec:ch4_feature_comparison},
где модели изначально обучались на таком представлении.
Сравниваются разные задачи и процедуры подготовки сети.

Изменение качества по всем проверенным группам показано
на рисунке~\ref{fig:ch4_channel_dropout}.
Он позволяет сопоставить не только общие токовые
и напряженческие ветви, но и отдельные гармоники.
Высота столбца характеризует снижение Macro-F1
при заданном занулении, а не вероятность отказа
измерительного канала в эксплуатации.

\begin{figure}[p]
 \centering
 \includegraphics[width=0.94\textwidth,height=0.78\textheight,keepaspectratio]{ch04_channel_dropout_article.png}
 \caption{Снижение Macro-F1 при последовательном
 занулении групп спектральных признаков}
 \label{fig:ch4_channel_dropout}
\end{figure}
\begin{figure}[p]
 \centering
 \includegraphics[width=0.94\textwidth,height=0.78\textheight,keepaspectratio]{ch04_channel_dropout_classes_article.png}
 \caption{Изменение F1 отдельных классов ОЗЗ
 при занулении тех же групп входа}
 \label{fig:ch4_channel_dropout_classes}
\end{figure}
\figurereview{Исходные рисунки 7 и 8 статьи 4. Перевести имена
групп входа и классов ОЗЗ, увеличить шрифт вертикальной легенды.
Для all\_angles писать «зануление углов», а не «потеря фаз».
Исходные данные и код рисунков не дублировались.}

Покомпонентное представление на
рисунке~\ref{fig:ch4_channel_dropout_classes}
показывает, что одно и то же изменение входа
по-разному влияет на четыре класса.
Небольшое среднее снижение не исключает
заметной потери для отдельного вида повреждения.
Поэтому общий Macro-F1 рассматривается вместе
с четырьмя частными оценками.

% Карты зануления относятся к предыдущему опыту.
\clearpage
\subsection{Градиентная чувствительность к сигналам и времени}
\label{sec:ch4_gradient_sensitivity}
\conceptsource{Статья 4, градиентный анализ и рисунки 9--10.
Объяснён смысл Gradient×Input, роль фазовых напряжений,
опорной фазы и внутренних позиций окна. Это не измерение
весов внимания и не доказанная задержка детекции.}

% D04-011 R4-16 BEGIN
\ifdefined\AuthorReview
\par\smallskip\begingroup\small\raggedright\color{red!60!black}
\noindent\textbf{Примечание рецензента [R4-16], редактура рисунков.} При просмотре PDF наиболее трудно читаются мелкие легенды RTDS-примеров, групп зануления и временной градиентной карты. Это вопрос доступности результата, а не совпадения шрифта. Для окончательной версии перевести сокращения и сделать читаемыми именно поля, на которые опирается вывод; часть однотипных полноформатных примеров можно перенести в приложение после согласования. Исходные файлы сохранять, формы кривых не дорисовывать.\par\endgroup\smallskip
\fi
% D04-011 R4-16 END
\ifdefined\AuthorReview
\par\smallskip\begingroup\small\raggedright\color{blue!55!black}
\noindent\textbf{Ответ на R4-16, D04-012-B.} Задание на укрупнение легенд и разделение временной карты добавлено в прежний figure\_specs/ch04\_physkan.md. Однотипные примеры не перенесены в приложение без согласования. Исходные изображения сохранены, кривые не дорисовывались; графическая правка ещё требуется.\par\endgroup\smallskip
\fi

(Примечание автора: Ок, по сути это задача для меня для правок, не удаляем это, пока не будет сделано.)

Программное отключение заметно меняет вход.
Для рассмотрения локальной чувствительности используется
также градиентный анализ. Вычисляется, как меняется
выходная оценка класса при малом изменении каждого
входного признака. Полученная производная умножается
на само значение признака. Для карты берётся модуль
этого произведения, после чего величины
объединяются по выбранным сигналам или времени.
Так строится карта Gradient×Input.

В отличие от весов самовнимания, эта карта относится
к чувствительности итогового выхода сети.
Она зависит от выбранного класса, примера,
нормировки и способа объединения величин.
Её численное значение не является электрической
мощностью или долей физического вклада канала.

На рисунке~\ref{fig:ch4_gradient_signals}
показано распределение по измерительным сигналам.
Для рассматриваемой модели высокие значения
наблюдаются у фазных напряжений,
особенно для динамичного класса Белякова.
Различие между фазами может быть связано
с составом примеров и выбором опорного напряжения,
относительно которого отсчитываются фазовые углы.
Оно не означает, что одна фаза физически
важнее остальных при любом ОЗЗ.

\begin{figure}[htbp]
 \centering
 \includegraphics[width=\textwidth]{ch04_gradient_signals_article.png}
 \caption{Объединённая градиентная чувствительность
 выходов ОЗЗ к измерительным сигналам}
 \label{fig:ch4_gradient_signals}
\end{figure}
\figurereview{Исходный рисунок 9 статьи 4. Перевести названия классов
и Mean|grad×input|; не переименовывать карту в «веса внимания».
Числа характеризуют выбранную нормировку градиентной оценки.}

Карта по временным положениям приведена
на рисунке~\ref{fig:ch4_gradient_time}.
Она показывает распределение чувствительности
внутри поданного окна и позволяет выделить
спектральные координаты с большим откликом.
Начальная позиция окна здесь является внутренним
индексом представления. Без дополнительной привязки
к моменту повреждения она не задаёт время
от начала аварии до решения органа.

\begin{figure}[p]
 \centering
 \includegraphics[width=\textwidth]{ch04_gradient_time_article.png}
 \caption{Градиентная чувствительность к спектральным
 признакам в разных положениях входного окна}
 \label{fig:ch4_gradient_time}
\end{figure}
\figurereview{Исходный рисунок 10 статьи 4. Перевести названия
каналов и Time step. Проверить отображение всех72 зон и масштаб
цвета отдельно для каждого класса; не трактовать самые левые
столбцы как автоматически доказанную реакцию в первые миллисекунды.}

Совместно две проверки дают более полное описание
используемого входа. Зануление показывает изменение
качества при крупном воздействии на группу,
а градиентная карта локальную чувствительность
для рассматриваемых примеров. Ни одна из них
по отдельности не восстанавливает полное
физическое правило решения нейросети.

\subsection{Ограничения проверки и дальнейшее использование}
\label{sec:ch4_input_sensitivity_meaning}
\conceptsource{Статья 4. Пределы имитации отказа и различие роли входов в детекции ОЗЗ и РНМ.}


Программное зануление группы не полностью воспроизводит
отказ измерительной цепи. При реальном нарушении
могут одновременно изменяться масштаб, полярность,
шум, насыщение и достоверность нескольких каналов.
Кроме того, входные преобразования могут рассчитывать
связанные признаки из тех же исходных измерений.
Поэтому для проверки реального отказа требуется
учесть всё распространение изменения по тракту.

Рассмотренные результаты дополняют выбор входа.
В задаче ОЗЗ они показывают роль фазных напряжений,
угловой информации и гармонического состава.
Они также указывают группы, для которых особенно
существенны контроль качества измерения
и согласование подготовки сигналов при обучении
и последующем применении.

Для исследования БАВР это основание сохранить
связанные токовые и напряженческие измерения
в общей подготовке данных и выбирать их использование
по назначению конкретного органа.
Слабое влияние группы в одной проверке
не является основанием удалить её из набора
для всех последующих задач.
